\subsection{分辨率折叠映射（folding）}\label{subsec:folding-map}

折叠的目标在于将任意窗口微态 $\omega=\omega_1\cdots \omega_m\in\Omega_m$ 规范化为一个合法词 $x\in X_m$。

为避免“人为编号”，本文直接以 Fibonacci 权重赋予微态可审计整数值，再取 Zeckendorf 规范形：
\begin{definition}[Fibonacci 加权值与折叠算子]\label{def:fold-word}
对 $\omega=\omega_1\cdots\omega_m\in\Omega_m$ 定义其 Fibonacci 加权值
\[
N(\omega):=\sum_{k=1}^{m}\omega_k F_{k+1}\ \in\ \NN.
\]
令 $Z(N(\omega))=(c_1,c_2,\dots)$ 为 $N(\omega)$ 的 Zeckendorf 数字（\cite{Zeckendorf1972}），定义分辨率 $m$ 的折叠
\[
\Fold_m:\Omega_m\to X_m,\qquad
\Fold_m(\omega):=\pi_m\!\bigl(Z(N(\omega))\bigr).
\]
\leanverified{Fold}
\end{definition}

\subsubsection{折叠的局部重写实现：终止、合流与顺序无关}\label{subsec:fold-local-rewrite}
定义 \ref{def:fold-word} 以“先取值、再取 Zeckendorf 规范形”给出 $\Fold_m$。为将其落实为\textbf{局部、可执行}的归一化过程，并澄清“归约顺序不影响正规形”，我们把输入窗口嵌入到允许中间态出现冗余数字的配置空间上实施值保持重写。

\paragraph{配置空间与值函数}
令
\[
\mathcal{D}:=\Bigl\{a=(a_1,a_2,\dots): a_k\in\NN,\ \text{仅有限个 }a_k\neq 0\Bigr\}
\]
为有限支撑的非负整数数字配置。定义值同态
\[
\mathrm{Val}(a):=\sum_{k\ge 1} a_k F_{k+1}.
\]
对给定 $\omega\in\Omega_m$，记其嵌入为
\[
\iota_m(\omega):=(\omega_1,\dots,\omega_m,0,0,\dots)\in\mathcal{D},
\]
则 $\mathrm{Val}(\iota_m(\omega))=N(\omega)$。

\paragraph{局部重写规则（值保持）}
对 $a\in\mathcal{D}$，允许在满足条件时对局部模式做如下替换（其中 $e_k$ 为第 $k$ 位单位向量）：
\begin{enumerate}
  \renewcommand{\theenumi}{R\arabic{enumi}}
  \renewcommand{\labelenumi}{(\theenumi)}
  \item \label{rule:fold:adj}\textbf{相邻合并：}若 $a_k\ge 1$ 且 $a_{k+1}\ge 1$，则
  \[
  a\ \to\ a-e_k-e_{k+1}+e_{k+2}.
  \]
  \item \label{rule:fold:dedup1}\textbf{底位去重：}若 $a_1\ge 2$，则
  \[
  a\ \to\ a-2e_1+e_2.
  \]
  \item \label{rule:fold:dedup2}\textbf{第二位去重：}若 $a_2\ge 2$，则
  \[
  a\ \to\ a-2e_2+e_1+e_3.
  \]
  \item \label{rule:fold:dedupk}\textbf{一般位去重（$k\ge 3$）：}若 $a_k\ge 2$，则
  \[
  a\ \to\ a-2e_k+e_{k-2}+e_{k+1}.
  \]
\end{enumerate}
每条规则都只改动有限个坐标，因此是严格的\textbf{局部}重写；并且它们分别对应恒等式
\(
F_{k+1}+F_{k+2}=F_{k+3},
2F_2=F_3,
2F_3=F_2+F_4,
2F_{k+1}=F_{k-1}+F_{k+2}
\)，因此每一步都保持 $\mathrm{Val}$ 不变。

\begin{lemma}[二元局部归一化的权重递推刚性]\label{lem:fold-local-weight-rigidity-fibonacci}
设同一组局部规则 \textup{(\ref{rule:fold:adj})--(\ref{rule:fold:dedupk})} 作用在
\(
\mathrm{Val}_w(a):=\sum_{k\ge 1}a_k w_k
\)
上并要求每一步重写都严格保持 \(\mathrm{Val}_w\) 不变。
则权重序列 \((w_k)_{k\ge 1}\) 必满足
\[
w_{k+2}=w_{k+1}+w_k\qquad(k\ge 1).
\]
若进一步固定归一化 \(w_1=1,\ w_2=2\)，则唯一解为 \(w_k=F_{k+1}\)。
\end{lemma}

\begin{proof}
由规则 \textup{(\ref{rule:fold:adj})} 的值保持，
\[
w_k+w_{k+1}=w_{k+2}\qquad(k\ge 1),
\]
即得二阶线性递推。
该递推已蕴含规则 \textup{(\ref{rule:fold:dedup1})--(\ref{rule:fold:dedupk})} 的值保持恒等式，因此为必要条件。
在初值 \(w_1=1,w_2=2\) 下，递推唯一确定 \(w_k=F_{k+1}\)。
\leanverified{weight\_rigidity\_fibonacci}
\leanverified{local\_weight\_rigidity\_step}
\leanverified{local\_weight\_rigidity\_normalized}
\leanverified{paper\_fold\_local\_weight\_rigidity\_fibonacci}
\leanverified{local\_weight\_rigidity\_first\_four}
\leanverified{fold\_rule\_R2\_identity}
\leanverified{fold\_rule\_R3\_identity}
\leanverified{fold\_rule\_R4\_identity}
\leanverified{paper\_fold\_local\_rewrite\_identities}
\leanverified{fold\_rule\_R4\_gap}
\leanverified{fold\_rule\_completeness\_seeds}
\end{proof}

\paragraph{不可约项与 Zeckendorf 规范形}
记 $X_\infty=\{(c_1,c_2,\dots)\in\{0,1\}^{\mathbb{N}}:\ c_kc_{k+1}=0\}$。上述规则恰在出现相邻 $11$ 或某位数字 $\ge 2$ 时可用，因此不可约项当且仅当属于 $X_\infty$。

\begin{proposition}[终止性、局部合流与合流（纽曼引理）]\label{prop:fold-rewrite-newman}
上述重写系统在 $\mathcal{D}$ 上强终止，且局部合流；由纽曼引理可得其合流性。因此每个 $a\in\mathcal{D}$ 的正规形存在且唯一，记为 $\mathrm{NF}(a)\in X_\infty$，并满足
\[
\mathrm{Val}(\mathrm{NF}(a))=\mathrm{Val}(a).
\]
\leanverified{step\_stronglyTerminating}
\leanverified{step\_confluent}
\leanverified{paper\_fold\_rewrite\_newman}
\end{proposition}

\begin{proof}
\textbf{（终止性）}令 $A(a)=\sum_{k\ge 1} a_k$、$B(a)=\sum_{k\ge 1} k\,a_k$、$C(a)=a_2$，并设 $\mu(a)=(A(a),B(a),C(a))$ 按字典序比较。则规则 (\ref{rule:fold:adj}) 与 (\ref{rule:fold:dedup1}) 使 $A$ 减少；规则 (\ref{rule:fold:dedupk})（$k\ge 3$）使 $B$ 减少；规则 (\ref{rule:fold:dedup2}) 保持 $A,B$ 不变但使 $C$ 减少。故每步重写都严格降低 $\mu$，系统强终止。

\textbf{（局部合流）}若 $a$ 有两条一步归约 $a\to a_1,a_2$，则两侧继续归约到不可约形后都落在 $X_\infty$，且由于每步保持 $\mathrm{Val}$，两不可约形同值；由 Zeckendorf 唯一性它们相同，于是两支可汇合。

\textbf{（合流）}由纽曼引理（强终止 + 局部合流 $\Rightarrow$ 合流）即得结论。
\end{proof}

\begin{corollary}[\texorpdfstring{$\Fold_m$}{Fold\_m} 的归约顺序无关实现]\label{cor:foldm-order-indep}
对任意 $\omega\in\Omega_m$，
\[
\Fold_m(\omega)\ =\ \pi_m\bigl(\mathrm{NF}(\iota_m(\omega))\bigr),
\]
且右端与归约时选取规则/位置的顺序无关。
\leanverified{normalPrefix\_iota\_eq\_Fold}
\leanverified{paper\_foldm\_order\_indep}
\end{corollary}

\begin{remark}[规范截面与可计算性]\label{rem:fold-normal-form-computable-section}
命题 \ref{prop:fold-rewrite-newman} 给出一个可计算的规范截面：对任意 $a\in\mathcal{D}$，正规形 $\mathrm{NF}(a)\in X_\infty$ 唯一存在且与归约顺序无关。因此，当“第三参数”仅对应于等价类中代表的选择（规范冗余）时，无需显式记录该参数；可通过求正规形在多项式时间内选出唯一代表（分辨率口径为线性，见命题 \ref{prop:fold-complexity}）。需要强调的是，该消去仅恢复稳定代表，不恢复同一纤维内的微态残差。
\end{remark}

\begin{proposition}[折叠的规范性、幂等性与满射性]\label{prop:fold-basic}
折叠算子 $\Fold_m$ 良定义且满足：
\begin{enumerate}
\normalfont
  \item \textbf{（规范性）} $\Fold_m(\omega)\in X_m$；
  \item \textbf{（幂等性）} 若 $\omega\in X_m$，则 $\Fold_m(\omega)=\omega$；
  \item \textbf{（满射性）} $\Fold_m:\Omega_m\twoheadrightarrow X_m$ 为满射。
\end{enumerate}
\leanverified{Fold\_stable}
\leanverified{Fold\_idempotent}
\leanverified{Fold\_surjective}
\end{proposition}
\begin{definition}[局部重写证书读出]\label{def:distill-stallings-local-rewrite-certified-readout}
令 \(S\) 为集合。称映射 \(R:\mathcal D\to S\) 为局部重写证书读出，若对每条允许的局部重写 \(a\to b\) 都有
\[
R(a)=R(b).
\]
固定 \(m\ge1\)。若 \(R\) 还满足：对任意 \(c,c'\in X_\infty\)，只要 \(\pi_m(c)=\pi_m(c')\) 就有 \(R(c)=R(c')\)，则称 \(R\) 为 \(m\)-核心读出。
\end{definition}

\begin{proposition}[局部证书读出只能经过折叠核心]\label{prop:distill-stallings-local-rewrite-certified-readout-factorization}
设 \(R:\mathcal D\to S\) 为任意映射。
\begin{enumerate}
  \item \(R\) 是局部重写证书读出，当且仅当存在唯一映射
  \[
  \widehat R:X_\infty\to S
  \]
  使得
  \[
  R=\widehat R\circ\mathrm{NF}.
  \]
  \item 固定 \(m\ge1\)。\(R\) 是 \(m\)-核心读出，当且仅当存在唯一映射
  \[
  \bar R:X_m\to S
  \]
  使得
  \[
  R=\bar R\circ\pi_m\circ\mathrm{NF}.
  \]
\end{enumerate}
因此，任何只登记局部重写保持性的读出，都不能把中间冗余数字、重写路径或可修剪尾部当作稳定对象；若还要求分辨率 \(m\) 可见，则它必须经由稳定前缀 \(\pi_m(\mathrm{NF}(-))\)。
\end{proposition}
\begin{proposition}[稳定核心商的可采纳性判别与最坏坏例]
\label{distill:stallings_foldings_and_bestvina_handel_train_tracks-local_fold_preserves_represented_object-stable-core-quotient-worst-witness}
固定 \(m\ge1\)，并记
\[
\kappa_m:\mathcal D\longrightarrow X_m,
\qquad
\kappa_m(a):=\pi_m(\mathrm{NF}(a)).
\]
令 \(E\subset\mathcal D\) 为非空有限集合，令 \(q:E\twoheadrightarrow Q\) 为一个有限商映射。以下条件等价：
\begin{enumerate}
  \item \(q\) 是稳定核心可采纳的，即
  \[
  q(a)=q(b)\quad\Longrightarrow\quad \kappa_m(a)=\kappa_m(b)
  \qquad(a,b\in E);
  \]
  \item 存在唯一映射 \(\bar\kappa:Q\to X_m\)，使
  \[
  \kappa_m|_E=\bar\kappa\circ q;
  \]
  \item 对任意集合 \(S\) 与任意 \(m\)-核心读出
  \[
  R:E\to S,
  \qquad
  R=\bar R\circ\kappa_m|_E
  \quad(\bar R:X_m\to S),
  \]
  存在唯一 \(\widetilde R:Q\to S\) 使
  \[
  R=\widetilde R\circ q;
  \]
  \item 对每个子集 \(A\subseteq X_m\)，谓词
  \[
  a\longmapsto \mathbf 1_A(\kappa_m(a))
  \]
  都可下降到 \(Q\)。
\end{enumerate}
若这些条件失败，则存在二点坏例
\[
 a,b\in E,
 \qquad q(a)=q(b),
 \qquad \kappa_m(a)\ne\kappa_m(b),
\]
并且单点稳定谓词
\[
A=\{\kappa_m(a)\}\subseteq X_m
\]
已经检测到该失败：\(\mathbf 1_A\circ\kappa_m\) 不能下降到 \(Q\)。

特别地，若 \(B\subseteq\Omega_m\) 且 \(E=\iota_m(B)\)，则上述可采纳性等价于
\[
q(\iota_m(\omega))=q(\iota_m(\omega'))
\quad\Longrightarrow\quad
\Fold_m(\omega)=\Fold_m(\omega') .
\]
因此，窗口层的分支识别保存的不是原始配置 \(\omega\)，而是经由
\(\Fold_m\) 的稳定核心读出。另一方面，任何由规则
\textup{(\ref{rule:fold:adj})--(\ref{rule:fold:dedupk})} 的有限重写链生成的商都自动满足上述可采纳性。
\leanverified{paper_stallings_foldings_local_fold_preserves_represented_object_stable_core_quotient_worst_witness}
\end{proposition}
\begin{theorem}[可见尾影响的局部支撑骨架]\label{distill:stallings_foldings_and_bestvina_handel_train_tracks-finite_descent_to_folded_core-visible-tail-causal-core}
沿用小节 \ref{subsec:fold-local-rewrite} 的局部重写规则族
\(\mathcal R_{\rm Fib}\)。对一步重写 \(u\to v\)，记
\[
\operatorname{Supp}(u,v):=\{j:u_j\ne v_j\}.
\]
令
\[
D_{\rm Fib}:=\max\{\max S-\min S:S=\operatorname{Supp}(u,v),\ u\to v\text{ 为允许一步重写}\}.
\]
则 \(D_{\rm Fib}\le 3\)。对 \(M\ge1\) 记
\[
a_{\le M}:=\sum_{1\le k\le M}a_k e_k .
\]
固定 \(q\le M\)，并令
\[
\tau:\quad a=a^{(0)}\to a^{(1)}\to\cdots\to a^{(L)}=\mathrm{NF}(a)
\]
为任一合法归一化链。构造一个有限无向图
\(\mathsf C_{q\mid M}(\tau,a)\)：其顶点由两个形式顶点
\(\star_R,\star_T\)、步骤顶点 \(s_1,\ldots,s_L\)，以及所有坐标顶点
\[
c_j\qquad j\in \operatorname{supp}(a)\cup\bigcup_{i=1}^L
\operatorname{Supp}(a^{(i-1)},a^{(i)})
\]
组成；若 \(j\in\operatorname{Supp}(a^{(i-1)},a^{(i)})\)，连边 \(c_j-s_i\)；若
\(j\le q\)，连边 \(\star_R-c_j\)；若 \(j>M\) 且 \(a_j>0\)，连边
\(\star_T-c_j\)。
若
\[
\pi_q\mathrm{NF}(a)\ne \pi_q\mathrm{NF}(a_{\le M}),
\]
则对每一条合法归一化链 \(\tau\)，顶点 \(\star_R\) 与 \(\star_T\) 在
\(\mathsf C_{q\mid M}(\tau,a)\) 中属于同一连通分支。并且任一连接二者的路径若经过
\(s\) 个步骤顶点，则
\[
s\ge \left\lceil \frac{M+1-q}{D_{\rm Fib}}\right\rceil .
\]
等价地，若某条归一化链的 \(\star_R\)-分支不含 \(\star_T\)，则 \(M\) 之外的初始尾部相对于读出
\(\pi_q\mathrm{NF}(-)\) 是可修剪的 transient tail，并有
\[
\pi_q\mathrm{NF}(a)=\pi_q\mathrm{NF}(a_{\le M}).
\]
这里的一步重写正是本文语境中的 admissible local fold；上述连通分支是相对于基读出
\(\pi_q\) 的 folded visible core，而其余分支是由该读出与修剪规则共同确定的非核心尾部。
\end{theorem}
\begin{proof}
各规则支撑的直径可直接读出：规则 \textup{(\ref{rule:fold:adj})} 与
\textup{(\ref{rule:fold:dedup2})} 的支撑包含在长度为 \(3\) 的坐标区间内，规则
\textup{(\ref{rule:fold:dedup1})} 的支撑直径为 \(1\)，规则
\textup{(\ref{rule:fold:dedupk})} 的支撑包含在
\([k-2,k+1]\) 中，故 \(D_{\rm Fib}\le3\)。命题
\ref{prop:fold-rewrite-newman} 给出终止性、合流性与正规形的唯一性，因此可在任意合法归一化链上讨论同一终态
\(\mathrm{NF}(a)\)。

先证连通性断言的逆否命题。设对某条链 \(\tau\)，\(\star_R\) 与 \(\star_T\) 不连通。令
\(K\) 为 \(\star_R\) 所在连通分支。若步骤顶点 \(s_i\in K\) 而 \(s_j\notin K\)，则
\[
\operatorname{Supp}(a^{(i-1)},a^{(i)})\cap
\operatorname{Supp}(a^{(j-1)},a^{(j)})=\varnothing,
\]
否则二者会经由某个坐标顶点相连。故 \(K\) 内步骤与 \(K\) 外步骤的作用支撑两两不交。每条局部规则的适用性与效果只由其支撑中的坐标决定，所以 \(K\) 内步骤可与 \(K\) 外步骤交换次序，并且其在 \(K\) 内坐标上的演化不受 \(K\) 外步骤影响。

又因 \(\star_T\notin K\)，凡 \(j>M\) 且 \(a_j>0\) 的初始尾坐标均不属于 \(K\)。于是从 \(a\) 改为
\(a_{\le M}\) 不改变 \(K\) 内所有坐标的初值。按 \(\tau\) 中出现的相对次序执行全部 \(K\) 内步骤，在输入 \(a\) 与输入
\(a_{\le M}\) 上得到完全相同的 \(K\) 内坐标演化；特别地，坐标 \(1,\ldots,q\) 的结果相同。

在完整链 \(\tau\) 中，\(K\) 外步骤不触及 \(K\)，故从该时刻到终态
\(\mathrm{NF}(a)\) 为止，坐标 \(1,\ldots,q\) 不再改变。另一方面，把初始尾部删去只会把 \(K\) 外的若干数字减小为零，不会在穿越
\(K\) 与其补集的边界处制造新的相邻 \(11\) 或新的重数 \(\ge2\)；若某条后续归约需要触及 \(K\)，则在原链终态中已有同一局部非法模式，矛盾于
\(\mathrm{NF}(a)\in X_\infty\)。因此 \(a_{\le M}\) 的任意归一化补全均不再改变坐标
\(1,\ldots,q\)。由合流性，该补全的终态为 \(\mathrm{NF}(a_{\le M})\)，从而
\[
\pi_q\mathrm{NF}(a)=\pi_q\mathrm{NF}(a_{\le M}).
\]
这证明了若两端读出不同，则每条归一化链中 \(\star_R\) 与 \(\star_T\) 必连通。

最后证明预算下界。取一条从 \(\star_R\) 到 \(\star_T\) 的简单路径，并按路径顺序列出其中出现的步骤顶点
\(s_{i_1},\ldots,s_{i_s}\)。路径从某个坐标 \(\le q\) 进入第一个步骤支撑，最后从某个初始尾坐标
\(\ge M+1\) 离开最后一个步骤支撑；相邻步骤支撑之间经由坐标顶点相连，因而相交。每个步骤支撑的直径至多为
\(D_{\rm Fib}\)。一串两两相继相交、直径至多为 \(D_{\rm Fib}\) 的支撑若从坐标
\(\le q\) 连接到坐标 \(\ge M+1\)，其步骤数 \(s\) 必满足
\[
sD_{\rm Fib}\ge M+1-q.
\]
取上整数即得所示下界。证毕。
\end{proof}

\begin{proof}
先证前四项等价。第一项正是说 \(\kappa_m|_E\) 在 \(q\) 的每条纤维上常值，因此可定义
\[
\bar\kappa(q(a)):=\kappa_m(a).
\]
该定义良定当且仅当第一项成立；唯一性来自 \(q\) 的满射性。这证明第一项与第二项等价。

若第二项成立，并且 \(R=\bar R\circ\kappa_m|_E\)，则
\[
R=\bar R\circ\bar\kappa\circ q,
\]
故可取 \(\widetilde R:=\bar R\circ\bar\kappa\)。唯一性仍由 \(q\) 满射给出。故第二项推出第三项。反过来，在第三项中取 \(S=X_m\) 与 \(R=\kappa_m|_E\)，便得到 \(\kappa_m|_E\) 经由 \(q\) 因子化，故第三项推出第二项。

第三项显然推出第四项，因为每个 \(\mathbf 1_A\circ\kappa_m\) 都是 \(m\)-核心读出。若第四项成立且 \(q(a)=q(b)\)，取
\(A=\{\kappa_m(a)\}\)。下降性给出
\[
1=\mathbf 1_A(\kappa_m(a))=\mathbf 1_A(\kappa_m(b)),
\]
于是 \(\kappa_m(b)=\kappa_m(a)\)。故第四项推出第一项。

若可采纳性失败，则按第一项的否定可取 \(a,b\in E\) 满足
\(q(a)=q(b)\) 但 \(\kappa_m(a)\ne\kappa_m(b)\)。令
\(A=\{\kappa_m(a)\}\)，则
\[
\mathbf 1_A(\kappa_m(a))=1,
\qquad
\mathbf 1_A(\kappa_m(b))=0.
\]
同一商纤维上取值不同，故该谓词不能下降到 \(Q\)。这给出所述二点坏例。

若 \(E=\iota_m(B)\)，由推论 \ref{cor:foldm-order-indep}，对每个 \(\omega\in\Omega_m\) 有
\[
\kappa_m(\iota_m(\omega))
=\pi_m(\mathrm{NF}(\iota_m(\omega)))
=\Fold_m(\omega),
\]
从而可采纳性等价于 \(\Fold_m\) 在商纤维上常值。

最后，若 \(a,b\) 由一条有限重写链连接，则命题 \ref{prop:fold-rewrite-newman} 的合流性给出
\(\mathrm{NF}(a)=\mathrm{NF}(b)\)，因而
\(\kappa_m(a)=\kappa_m(b)\)。所以由这些局部重写链生成的商只会识别同一稳定核心纤维内的点，自动满足第一项。证毕。
\end{proof}

\begin{proof}
先证第一项。若 \(R\) 对每条局部重写不变，则由命题 \ref{prop:fold-rewrite-newman}，任意 \(a\in\mathcal D\) 有有限归约链
\[
a=a^{(0)}\to a^{(1)}\to\cdots\to a^{(N)}=\mathrm{NF}(a).
\]
沿链逐步使用不变性，得
\[
R(a)=R(\mathrm{NF}(a)).
\]
于是可定义 \(\widehat R(c):=R(c)\) 对 \(c\in X_\infty\)，并有 \(R=\widehat R\circ\mathrm{NF}\)。唯一性来自 \(\mathrm{NF}|_{X_\infty}=\mathrm{id}\)。反之，若 \(R=\widehat R\circ\mathrm{NF}\)，且 \(a\to b\)，则合流性给出 \(\mathrm{NF}(a)=\mathrm{NF}(b)\)，故 \(R(a)=R(b)\)。

第二项中，若 \(R\) 为 \(m\)-核心读出，则第一项给出 \(R=\widehat R\circ\mathrm{NF}\)。定义
\[
\bar R(x):=\widehat R(c),\qquad c\in X_\infty,\ \pi_m(c)=x.
\]
这是良定义的，因为 \(m\)-核心条件正是说 \(\widehat R\) 在 \(\pi_m\)-纤维上常值。于是 \(R=\bar R\circ\pi_m\circ\mathrm{NF}\)，且唯一性由 \(\pi_m:X_\infty\to X_m\) 的满射性给出。反向显然：任何经由 \(\pi_m\circ\mathrm{NF}\) 的读出既保持局部重写，又在不可约项的同一 \(m\)-前缀纤维上常值。证毕。
\end{proof}

\begin{proof}
\textbf{（规范性）}由定义 \ref{def:fold-word}，$Z(N(\omega))\in X_\infty$，故 $\pi_m(Z(N(\omega)))\in X_m$。

\textbf{（幂等性）}若 $\omega\in X_m$，令 $N=\sum_{k=1}^m \omega_k F_{k+1}$。由于 $\omega$ 已是合法 Zeckendorf 数字的长度-$m$ 前缀，且 Zeckendorf 表示唯一，故 $Z(N)$ 的前 $m$ 位必等于 $\omega$，从而 $\Fold_m(\omega)=\omega$。

\textbf{（满射性）}对任意 $x\in X_m$，取 $\omega=x\in\Omega_m$，由幂等性得 $\Fold_m(\omega)=x$，故 $\Fold_m$ 满射。
\end{proof}
